\section{Introduction}
\label{sec:intro}

Diffusion-based structure generators such as AlphaFold~3 \citep{alphafold3},
Boltz-2 \citep{boltz2}, and RoseTTAFold All-Atom \citep{rfaa} are now used at
large scale, yet systematic evaluations continue to expose failure modes:
confidence can be miscalibrated on disordered regions \citep{gopalan2025},
protein--ligand co-folding can track memorized poses \citep{skrinjar2025}, and
antibody--antigen docking remains brittle at a single seed
\citep{hitawala2025}. A complementary diagnostic question is not only whether
a \emph{single} prediction looks plausible, but what distribution of outcomes a
model assigns to a \emph{fixed} input when its internal stochasticity is
exhausted. One forward pass answers the former; the latter requires sampling over
entire denoising trajectories under the generator's transition law.

Transition path sampling (TPS), developed for rare transitions in molecular
simulation \citep{dellago1998,bolhuis2002}, has recently been extended to
non-equilibrium and learned generators
\citep{buijsman2020,falkner2023,falkner2024}. The same construction applies when
each reverse-diffusion step is a deterministic map of explicit noise and
auxiliary random variables: the object of interest is the induced distribution
on reactive paths connecting a high-noise set $A$ and a structured endpoint
set $B$. Path-space Metropolis--Hastings moves then yield unbiased reactive
ensembles with computable log-weights. Parallel work applies path sampling ideas
to large language models \citep{dorman2026} and to diffusion-driven molecular
pathways via action minimization \citep{raja2025omtps}.

This paper develops the discrete-time path measure for the Boltz-2 reverse
sampler, including SE(3) augmentation and EDM-style churn \citep{karras2022},
gives the forward- and backward-shooting acceptance corrections following
\citet{buijsman2020,falkner2024}, analyzes Jacobian factors as diagnostics versus
as part of extended-variable weights, and combines TPS with on-the-fly
probability enhanced sampling (OPES) \citep{invernizzi2020,invernizzi2022} on
collective variables of the terminal frame. We relate rigid-body augmentation
to quotient-space and equivariant-network perspectives
\citep{xu2026quotient,yim2023se}. Nothing here claims an explicit likelihood for
the neural denoiser: the path weight uses Gaussian priors on injected noise and
augmentation together with the deterministic map from those variables to
coordinates; the network enters only through that map.

\paragraph{Contributions.}
\begin{enumerate}
  \item We state a generic discrete-time path distribution for a
    non-reversible generator and the Metropolis--Hastings acceptance for forward
    shooting, backward shooting with a dedicated backward proposal, and global
    reshuffle, including a simple irreducibility argument when reshuffle has
    positive weight.
  \item We specialize the construction to Boltz-2's extended state (rotations,
    translations, churn noise, coordinates), give the log path density used in
    code, and clarify when schedule-only Jacobian surrogates cancel in forward
    shooting and when they must not be inserted into backward proposals.
  \item We derive the multiplicative Boltzmann correction for OPES bias on
    terminal-frame collective variables and record caveats when the bias
    evolves during the chain.
  \item We report experiments on a 228-residue protein--ligand co-folding
    system: unbiased TPS shows sharply concentrated terminal structures, while
    TPS+OPES exposes metastable pose basins in a two-dimensional collective
    variable projection.
\end{enumerate}

\paragraph{Road map.}
Section~\ref{sec:related} situates the work in the literature.
Section~\ref{sec:tps-foundations} develops TPS on discrete paths.
Section~\ref{sec:diffusion-boltz} connects score-based diffusion and EDM
preconditioning to the discrete kernel relevant for sampling.
Section~\ref{sec:kernel-full} gives the Boltz-2 extended kernel, Jacobians, and
implementation log-ratios. Section~\ref{sec:opes-theory} treats OPES and biased
TPS. Section~\ref{sec:quotient} discusses quotient-space diffusion and
equivariance. Section~\ref{sec:experiments} presents experiments.
Sections~\ref{sec:discussion} and~\ref{sec:conclusion} discuss limitations and
close. With the problem and contributions fixed, the next section reviews prior
art and notation common to rare-event methods and generative modeling.

\paragraph{Notation.}
We write $x_i \in \R^{M \times 3}$ for the coordinate tensor at reverse-diffusion
index $i$ after padding to $M$ atoms. Steps run from $i=0$ (high noise) to
$i=L$ (structured output). We use $\mathcal{S}_i$ for the collection of random
objects fed into the sampler at step $i$. Inner products and norms are
Frobenius unless stated otherwise. We write $k_BT$ in OPES formulas for
compatibility with molecular simulation convention; in the pure TPS chain
without a thermal bath it should be read as an arbitrary energy scale fixing
the bias strength.
